\section{Motivation}

Modern data platforms increasingly need to react to events as they happen rather than waiting for the next batch load. A music-streaming service, for example, wants to know within seconds which songs are trending, whether a user's session just dropped, or whether ingestion is falling behind -- questions a traditional overnight ETL job cannot answer. Building and operating such a platform correctly requires more than a single tool: it requires a message queue that can buffer and replay events, a stream-processing engine that can keep up with continuous arrival, a storage layer that supports transactional updates on top of cheap object storage, a transformation layer that turns raw events into trustworthy analytics tables, and a way to observe the whole system's health while it runs.

This project, \emph{Streamify}, was motivated by the goal of building and demonstrating exactly such an end-to-end pipeline, using synthetic but realistic music-streaming events generated by Eventsim, rather than studying any single one of these components in isolation. Each stage of the pipeline -- ingestion, stream processing, lakehouse storage, transformation, and visualization -- was chosen and configured to work together as a single reproducible system, so that the trade-offs made at each stage (Chapter~\ref{ch:related-work-background}) could be evaluated in the context of the whole pipeline rather than as isolated technology choices.

A further motivation was to explore these technologies under realistic resource constraints. Rather than assuming unlimited cloud compute, the project deliberately targets a small, single-machine deployment (Section~\ref{sec:proposed-cluster-setup}), so that decisions such as Kafka cluster sizing and Spark cluster sizing (Section~\ref{sec:proposed-cluster-setup}) had to be made under the same resource pressure a small team or a proof-of-concept deployment would face in practice.
